\documentclass[10pt,a4paper]{amsart}
\usepackage{amsmath,amssymb,amsthm,graphicx}
% Retained minimal preamble for the Stacks Project chapter more-algebra.tex (More on Algebra).
% The chapter file pinned-original-source.tex carries only \input{preamble} on its first line, so no
% preamble text lives inside the chapter itself.  The Stacks Project uses one single shared preamble for
% every chapter; the snapshot stages it identically as sources/stacks-chow/preamble.tex and
% sources/stacks-etale-cohomology/preamble.tex.  The block below is copied mechanically, line for line,
% from sources/stacks-chow/preamble.tex lines 14-18 (the Xy-pic setup) and 154-217 (the native theorem
% declarations together with the mathematical macros).  Nothing was added, removed or redefined; the
% wrapper class and the amsmath/amssymb/amsthm/graphicx base packages are supplied by the assembler.
% Chapter: more-algebra.tex -- More on Algebra.
% For commutative diagrams we use Xy-pic
\usepackage[all]{xy}

% We use 2cell for 2-commutative diagrams.
\xyoption{2cell}

% Theorem environments.
%
\theoremstyle{plain}
\newtheorem{theorem}[subsection]{Theorem}
\newtheorem{proposition}[subsection]{Proposition}
\newtheorem{lemma}[subsection]{Lemma}

\theoremstyle{definition}
\newtheorem{definition}[subsection]{Definition}
\newtheorem{example}[subsection]{Example}
\newtheorem{exercise}[subsection]{Exercise}
\newtheorem{situation}[subsection]{Situation}

\theoremstyle{remark}
\newtheorem{remark}[subsection]{Remark}
\newtheorem{remarks}[subsection]{Remarks}

\numberwithin{equation}{subsection}

% Macros
%
\def\lim{\mathop{\mathrm{lim}}\nolimits}
\def\colim{\mathop{\mathrm{colim}}\nolimits}
\def\Spec{\mathop{\mathrm{Spec}}}
\def\Hom{\mathop{\mathrm{Hom}}\nolimits}
\def\Ext{\mathop{\mathrm{Ext}}\nolimits}
\def\SheafHom{\mathop{\mathcal{H}\!\mathit{om}}\nolimits}
\def\SheafExt{\mathop{\mathcal{E}\!\mathit{xt}}\nolimits}
\def\Sch{\mathit{Sch}}
\def\Mor{\mathop{\mathrm{Mor}}\nolimits}
\def\Ob{\mathop{\mathrm{Ob}}\nolimits}
\def\Sh{\mathop{\mathit{Sh}}\nolimits}
\def\NL{\mathop{N\!L}\nolimits}
\def\CH{\mathop{\mathrm{CH}}\nolimits}
\def\proetale{{pro\text{-}\acute{e}tale}}
\def\etale{{\acute{e}tale}}
\def\QCoh{\mathit{QCoh}}
\def\Ker{\mathop{\mathrm{Ker}}}
\def\Im{\mathop{\mathrm{Im}}}
\def\Coker{\mathop{\mathrm{Coker}}}
\def\Coim{\mathop{\mathrm{Coim}}}

% Boxtimes
%
\DeclareMathSymbol{\boxtimes}{\mathbin}{AMSa}{"02}

%
% Macros for moduli stacks/spaces
%
\def\QCohstack{\mathcal{QC}\!\mathit{oh}}
\def\Cohstack{\mathcal{C}\!\mathit{oh}}
\def\Spacesstack{\mathcal{S}\!\mathit{paces}}
\def\Quotfunctor{\mathrm{Quot}}
\def\Hilbfunctor{\mathrm{Hilb}}
\def\Curvesstack{\mathcal{C}\!\mathit{urves}}
\def\Polarizedstack{\mathcal{P}\!\mathit{olarized}}
\def\Complexesstack{\mathcal{C}\!\mathit{omplexes}}
% \Pic is the operator that assigns to X its picard group, usage \Pic(X)
% \Picardstack_{X/B} denotes the Picard stack of X over B
% \Picardfunctor_{X/B} denotes the Picard functor of X over B
\def\Pic{\mathop{\mathrm{Pic}}\nolimits}
\def\Picardstack{\mathcal{P}\!\mathit{ic}}
\def\Picardfunctor{\mathrm{Pic}}
\def\Deformationcategory{\mathcal{D}\!\mathit{ef}}
\usepackage{cleveref}
\crefname{originallemma}{original lemma}{original lemmas}
\Crefname{originallemma}{Original lemma}{Original lemmas}
\makeatletter
\newlabel{categories-definition-dual}{{0FFP}{source}{}{sourcepointer-1}{}}
\newlabel{categories-definition-dual@cref}{{[originallemma][0][]0FFP}{[1][0][]source}{}{}{}}
\makeatother
\crefname{originallemma}{original lemma}{original lemmas}
\Crefname{originallemma}{Original lemma}{Original lemmas}
\makeatletter
\newlabel{derived-lemma-subcategory-left-resolution}{{05T7}{source}{}{sourcepointer-2}{}}
\newlabel{derived-lemma-subcategory-left-resolution@cref}{{[originallemma][0][]05T7}{[1][0][]source}{}{}{}}
\makeatother
\crefname{originallemma}{original lemma}{original lemmas}
\Crefname{originallemma}{Original lemma}{Original lemmas}
\makeatletter
\newlabel{lemma-K-flat-resolution}{{14380}{source}{}{sourcepointer-3}{}}
\newlabel{lemma-K-flat-resolution@cref}{{[originallemma][0][]14380}{[1][0][]source}{}{}{}}
\makeatother
\crefname{originallemma}{original lemma}{original lemmas}
\Crefname{originallemma}{Original lemma}{Original lemmas}
\makeatletter
\newlabel{lemma-complex-bounded-above-free-colim-bounded-finite-free}{{38553}{source}{}{sourcepointer-4}{}}
\newlabel{lemma-complex-bounded-above-free-colim-bounded-finite-free@cref}{{[originallemma][0][]38553}{[1][0][]source}{}{}{}}
\makeatother
\crefname{originallemma}{original lemma}{original lemmas}
\Crefname{originallemma}{Original lemma}{Original lemmas}
\makeatletter
\newlabel{lemma-dual-perfect-complex}{{19176}{source}{}{sourcepointer-5}{}}
\newlabel{lemma-dual-perfect-complex@cref}{{[originallemma][0][]19176}{[1][0][]source}{}{}{}}
\makeatother
\crefname{originallemma}{original lemma}{original lemmas}
\Crefname{originallemma}{Original lemma}{Original lemmas}
\makeatletter
\newlabel{lemma-left-dual-complex}{{18462}{source}{}{sourcepointer-6}{}}
\newlabel{lemma-left-dual-complex@cref}{{[originallemma][0][]18462}{[1][0][]source}{}{}{}}
\makeatother
\crefname{originallemma}{original lemma}{original lemmas}
\Crefname{originallemma}{Original lemma}{Original lemmas}
\makeatletter
\newlabel{lemma-summands-perfect}{{18963}{source}{}{sourcepointer-7}{}}
\newlabel{lemma-summands-perfect@cref}{{[originallemma][0][]18963}{[1][0][]source}{}{}{}}
\makeatother
\title{An object of the derived category has a left dual precisely when it is perfect, and the dual is the dual of a perfect complex}
\author{The Stacks Project Authors}
\date{stacks-project:more-algebra.tex:local-snapshot-20261012; GNU FDL 1.2 or later}
\begin{document}
\maketitle
\setcounter{subsection}{1}
\noindent\textit{Editorial note.} All delivered mathematical text is the author's own, copied from the pinned Stacks Project chapter file more-algebra.tex; no statement or proof step was rewritten, reordered or removed.\par\medskip
\noindent\textit{Editorial note.} The chapter file pinned-original-source.tex carries only \textbackslash{}input\{preamble\} on its first line, so no preamble text lives inside the chapter. The retained file sources/stacks-more-algebra/retained-preamble.tex is a mechanical, line-for-line copy of the project's single shared preamble as staged in the snapshot; nothing was added, removed or redefined and no package was added.\par\medskip
\noindent\textit{Editorial note.} The source URL is the source-level Stacks Project URL, as required when more than one card is produced from one source; the per-claim tag is carried by claim\_key.\par\medskip
\noindent\textit{Editorial note.} The author's typo 'Choose a K-flat complexes N\textasciicircum{}bullet' is reproduced unchanged. The references inside the counted proof span are recorded as non-load-bearing navigation pointers: categories-definition-dual and derived-lemma-subcategory-left-resolution resolve to the project's cross-chapter tag list, while lemma-dual-perfect-complex, lemma-left-dual-complex, lemma-K-flat-resolution, lemma-complex-bounded-above-free-colim-bounded-finite-free and lemma-summands-perfect are same-chapter labels resolved to their \textbackslash{}label lines in the pinned source.\par\medskip
\noindent\textit{Allowed standard prerequisites (editorial).}\begin{itemize}
\item Left duals in a monoidal category and their unit and counit. the definition invoked at statement lines 38585-38586 and at proof line 38602. Supplies the maps eta and epsilon whose identities drive both factoring arguments.
\item K-flat complexes with flat terms and the identification of the derived tensor product with a total complex. used at proof lines 38605-38621. Makes Tot(M\textasciicircum{}bullet tensor\_R N\textasciicircum{}bullet) a representative of M tensor\textasciicircum{}L N and lets a finite description of eta(1) be extracted.
\item Bounded above complexes of free modules and bounded complexes of finite free modules. used at proof lines 38641-38650. Provides the re-representation of M that the second factoring argument cuts down to a perfect subcomplex.
\item Direct summands of perfect complexes are perfect. invoked at proof lines 38652-38653. Upgrades the factorisation of the identity through a perfect K to perfectness of M.
\end{itemize}\medskip
\ifdefined\hypertarget\hypertarget{sourcepointer-1}{}\fi
\noindent\textit{Source reference.} Original lemma 0FFP: focused-seeds/AGP/sources/stacks-snapshot-20261011/tags.txt, source lines 19294--19294. Cross-chapter navigation pointer into the Stacks Project; the standard fact it supplies is applied inside the counted proof span and the argument resting on it stays there.\par\medskip
\ifdefined\hypertarget\hypertarget{sourcepointer-2}{}\fi
\noindent\textit{Source reference.} Original lemma 05T7: focused-seeds/AGP/sources/stacks-snapshot-20261011/tags.txt, source lines 7471--7471. Cross-chapter navigation pointer into the Stacks Project; the standard fact it supplies is applied inside the counted proof span and the argument resting on it stays there.\par\medskip
\ifdefined\hypertarget\hypertarget{sourcepointer-3}{}\fi
\noindent\textit{Source reference.} Original lemma 14380: sources/stacks-more-algebra/pinned-original-source.tex, source lines 14380--14380. Same-chapter navigation pointer; the statement it names is standard for the section and the argument that consumes it stays inside the counted proof span.\par\medskip
\ifdefined\hypertarget\hypertarget{sourcepointer-4}{}\fi
\noindent\textit{Source reference.} Original lemma 38553: sources/stacks-more-algebra/pinned-original-source.tex, source lines 38553--38553. Same-chapter navigation pointer; the statement it names is standard for the section and the argument that consumes it stays inside the counted proof span.\par\medskip
\ifdefined\hypertarget\hypertarget{sourcepointer-5}{}\fi
\noindent\textit{Source reference.} Original lemma 19176: sources/stacks-more-algebra/pinned-original-source.tex, source lines 19176--19176. Same-chapter navigation pointer; the statement it names is standard for the section and the argument that consumes it stays inside the counted proof span.\par\medskip
\ifdefined\hypertarget\hypertarget{sourcepointer-6}{}\fi
\noindent\textit{Source reference.} Original lemma 18462: sources/stacks-more-algebra/pinned-original-source.tex, source lines 18462--18462. Same-chapter navigation pointer; the statement it names is standard for the section and the argument that consumes it stays inside the counted proof span.\par\medskip
\ifdefined\hypertarget\hypertarget{sourcepointer-7}{}\fi
\noindent\textit{Source reference.} Original lemma 18963: sources/stacks-more-algebra/pinned-original-source.tex, source lines 18963--18963. Same-chapter navigation pointer; the statement it names is standard for the section and the argument that consumes it stays inside the counted proof span.\par\medskip
\begin{lemma}
\label{lemma-have-dual-derived}
Let $R$ be a ring. Let $M$ be an object of $D(R)$. The following
are equivalent
\begin{enumerate}
\item $M$ has a left dual in $D(R)$ as in
Categories, Definition \ref{categories-definition-dual},
\item $M$ is a perfect object of $D(R)$.
\end{enumerate}
Moreover, in this case the left dual of $M$ is the object
$M^\vee$ of Lemma \ref{lemma-dual-perfect-complex}.
\end{lemma}

\begin{proof}
If $M$ is perfect, then we can represent $M$ by a bounded complex
$M^\bullet$ of finite projective $R$-modules. In this case $M^\bullet$
has a left dual in the category of complexes by
Lemma \ref{lemma-left-dual-complex}
which is a fortiori a left dual in $D(R)$.

\medskip\noindent
Assume (1). Say $N$, $\eta : R \to M \otimes_R^\mathbf{L} N$, and
$\epsilon : M \otimes_R^\mathbf{L} N \to R$ is a left dual
as in Categories, Definition \ref{categories-definition-dual}.
Choose a complex $M^\bullet$ representing $M$. Choose a
K-flat complexes $N^\bullet$ with flat terms
representing $N$, see Lemma \ref{lemma-K-flat-resolution}.
Then $\eta$ is given by a map of complexes
$$
\eta : R \longrightarrow \text{Tot}(M^\bullet \otimes_R N^\bullet)
$$
We can write the image of $1$ as a finite sum
$$
\eta(1) = \sum\nolimits_n \sum\nolimits_i m_{n, i} \otimes n_{-n, i}
$$
with $m_{n, i} \in M^n$ and $n_{-n, i} \in N^{-n}$. Let
$K^\bullet \subset M^\bullet$ be the subcomplex
generated by all the elements $m_{n, i}$ and $\text{d}(m_{n, i})$.
By our choice of $N^\bullet$ we find that
$\text{Tot}(K^\bullet \otimes_R N^\bullet) \subset
\text{Tot}(M^\bullet \otimes_R N^\bullet)$ and
$\eta(1)$ is in the subcomplex by our choice above.
Denote $K$ the object of $D(R)$ represented by $K^\bullet$.
Then we see that $\eta$ factors over a map
$\tilde \eta : R \longrightarrow K \otimes_R^\mathbf{L} N$.
Since
$(1 \otimes \epsilon) \circ (\eta \otimes 1) = \text{id}_M$
we conclude that the identity on $M$ factors through $K$
by the commutative diagram
$$
\xymatrix{
M \ar[rr]_-{\eta \otimes 1} \ar[rrd]_{\tilde \eta \otimes 1} & &
M \otimes_R^\mathbf{L} N \otimes_R^\mathbf{L} M
\ar[r]_-{1 \otimes \epsilon} &
M \\
& &
K \otimes_R^\mathbf{L} N \otimes_R^\mathbf{L} M
\ar[u] \ar[r]^-{1 \otimes \epsilon} &
K \ar[u]
}
$$
Since $K$ is bounded above it follows that $M \in D^-(R)$.
Thus we can represent $M$ by a bounded above complex
$M^\bullet$ of free $R$-modules, see for example
Derived Categories, Lemma \ref{derived-lemma-subcategory-left-resolution}.
Write $\eta(1) = \sum\nolimits_n \sum\nolimits_i m_{n, i} \otimes n_{-n, i}$
as before.
By Lemma \ref{lemma-complex-bounded-above-free-colim-bounded-finite-free}
we can find a subcomplex $K^\bullet \subset M^\bullet$
containing all the elements $m_{n, i}$
which is bounded and consists of finite free $R$-modules.
As above we find that the identity on $M$ factors through $K$.
Since $K$ is perfect we conclude $M$ is perfect too, see
Lemma \ref{lemma-summands-perfect}.
\end{proof}

\end{document}
